\documentclass[11pt,a4paper]{article}
\usepackage[margin=1in]{geometry}
\usepackage{amsmath,amssymb}
\usepackage{longtable}
\usepackage{booktabs}
\usepackage[hidelinks]{hyperref}
\usepackage[round]{natbib}
\usepackage{parskip}

\title{Literature position for the threshold-entry contest\\
\large Survey and verified bibliography}
\author{}
\date{1 September 2026}

\begin{document}
\maketitle

\section{Purpose and verification standard}

This document settles the literature check flagged as the gating open item in
\texttt{PROOFS.tex} (\S Open items, item 4) and in the referee handover
(R1 Comment 4). The question it answers is narrow: is the mechanism claimed
as novel --- an identical nominal entry cost $c$ made asymmetric in utility by
concavity, so that \emph{who} enters a fixed-prize contest is determined by the
wealth distribution rather than by ability --- already occupied?

Every entry in \texttt{refs.bib} has author list, title, journal, year, volume
and pages verified against at least one of: the article's own title page; the
publisher's article page; or the reference list of a published peer-reviewed
article. Where a claim below describes what a paper \emph{does} rather than what
it is, the evidence level is marked:

\begin{itemize}
\item[\textbf{[F]}] full text read;
\item[\textbf{[A]}] abstract or publisher summary only --- model details are
  reported as the abstract states them and have not been checked against the
  body;
\item[\textbf{[C]}] characterised via a citing paper's description, not the
  source itself.
\end{itemize}

Where the text read was a working paper or draft rather than the published
article, the tag says so --- ``[F, WP version]'' or ``[F, draft version]'' ---
and proposition numbers quoted from it may not match the published version.

Claims marked [A] or [C] are not safe to cite as statements about what a paper
proves. They are safe as statements about what a paper is about. Anything that
enters \texttt{PROOFS.tex} as a positioning claim should be upgraded to [F]
first.

\section{The closest paper on the affordability channel, read in full}
\label{sec:st}

\citet{SchroyenTreich2016} is the nearest occupant of the mechanism, and the
one against which the contribution has to be stated. It has now been read in
full \textbf{[F]}. The version read is TSE Working Paper 16-699 dated 11 April
2016, which the published \emph{Games and Economic Behavior} article replaced;
the published text has not been checked line by line against it, so any
quotation of a theorem number should be confirmed there before submission.

\paragraph{Correction to the previous draft of this document.} An earlier
version of this survey, written from the abstract, stated that their setting
produces no comparative statics on the wealth distribution. That is wrong.
Their Theorem~3 signs the effect of a mean-preserving spread in wealth in the
privilege contest. The overlap with P9 is therefore real and is set out below.

\paragraph{Primitives.} A static, one-shot, two-player game. Appendix~A.1 states
that the only source of heterogeneity is wealth $w_i$. Three models are
distinguished by whether the rent $r$ and the effort $x_i$ are commensurable
with wealth inside $u$:
\begin{align*}
  \text{privilege:} \quad & U_i = u(w_i - x_i) + \pi_i r, \\
  \text{ability:} \quad & U_i = \pi_i u(w_i + r) + (1-\pi_i) u(w_i) - c(x_i), \\
  \text{rent-seeking:} \quad & U_i = \pi_i u(w_i + r - x_i) + (1-\pi_i) u(w_i - x_i),
\end{align*}
with $\pi$ the power-logistic Tullock CSF of decisiveness $m$. Effort is
continuous on a convex feasible set. Dynamic effects are listed in their
conclusion as an extension not undertaken, as is an arbitrary number of players.

\paragraph{Consequence 1: the wealth/income distinction does no work.} Because
the model is static with no saving, $w_i$ is a one-period resource level out of
which effort is paid. That is formally the same object as the resource variable
in \texttt{PROOFS.tex}. Describing our primitive as income and theirs as wealth
is a labelling difference between two papers, not a difference between two
models, and it cannot be used to separate them. The separation has to come from
structure.

\paragraph{Consequence 2: our model is a privilege contest, confirmed from the
functional form.} Entry cost $c$ is paid out of resources and $V$ is separable
and exogenous, so only the marginal-cost channel is live and the
marginal-benefit channel is switched off by construction. This is now verified
against $U_i = u(w_i - x_i) + \pi_i r$ rather than inferred from the
taxonomy's name. It also fixes the open design question from the rebuild
handover: making $V$ money added to the resource would move the model to their
rent-seeking contest, where the two wealth effects oppose and, under CARA,
exactly cancel.

\paragraph{Consequence 3: the P9 overlap, stated honestly.} Their Theorem~3
gives, for the symmetric privilege contest, that a small mean-preserving spread
in wealth raises aggregate effort if and only if
\[
  2A(1-m^2) > P, \qquad A = -\frac{u''}{u'}, \qquad P = -\frac{u'''}{u''},
\]
reducing to $m < 2^{-1/2} \approx 0.707$ under CARA, to $m < 1$ under quadratic
utility, and to $\gamma(\tfrac12 - m^2) > \tfrac12$ under CRRA with coefficient
$\gamma$. So a mean-preserving spread in the resource distribution, inside a
privilege contest, is already in print.

\paragraph{Consequence 4: what survives, and why it is the contribution.} Their
result is local and heavily conditioned. It is a second-order effect of a small
spread evaluated at a \emph{symmetric} two-player equilibrium, obtained as a
quadratic form in the Hessian of the best-response function (their Theorem~2),
and its sign turns on the CSF decisiveness parameter and on the third derivative
of $u$. Their own summary concedes that the wealth-inequality effects are the
most complex in the paper and that signing a mean-preserving spread is very
difficult without further restricting functional forms. The condition
$u \in \Omega$ they need in the rent-seeking case --- that a mean-preserving
spread of background risk raises absolute risk aversion --- is, by
\citet{EeckhoudtGollierSchlesinger1996}, a restriction on higher-order risk
attitudes, and is stronger than DARA, being related to the risk vulnerability of
\citet{GollierPratt1996}.

P9-gen requires none of this. Its sign rule needs only that $\kappa$ is strictly
decreasing and that $T$ is a pivot-spread, with the flip located at
$w_{(k^*)}$ against the pivot $x_0$. No third derivative, no prudence condition,
no decisiveness parameter, no functional-form restriction, and the statement is
global rather than a local second-order expansion. The contribution sentence
should therefore be:

\begin{quote}
Moving the choice from the intensive to the extensive margin makes the
mean-preserving-spread comparative static sign-determinate without higher-order
restrictions on utility.
\end{quote}

That is a sharp claim stated directly against a named published result, and it
is stronger than anything the mechanism alone supports.

\paragraph{Consequence 5: a caveat that narrows the noise argument.} They note,
citing \citet{Jia2008}, that the power-logistic Tullock CSF can be derived as a
rank-order tournament in which the noise terms are inverse-exponential, and
remark that their contests may therefore also be read as rank-order
tournaments. The deterministic-versus-noisy contrast drawn against
\citet{HopkinsKornienko2010} in \S\ref{sec:hk} therefore does \emph{not} extend
to Schroyen--Treich. The remaining difference is narrower but real: their noise
is one specific family, whereas $F$ and $G$ here are arbitrary subject only to
the stated regularity.

\paragraph{Consequence 6: no entry margin.} There is no participation decision
anywhere in the paper. The zero-effort discontinuity in the best-response
functions is a corner of the continuous effort choice, not a fixed-cost entry
decision. The extensive margin is genuinely absent, which is what P5 and P7
require and what makes the contribution sentence above available.

\paragraph{A connection worth checking.} They contrast their
mean-preserving-spread results with the neutrality result of
\citet{BergstromBlumeVarian1986} for private provision of public goods, and note
in a footnote to their Theorem~2 that the quadratic form vanishes under the
Bergstrom--Varian conditions under which the distribution of agent
characteristics has no effect on aggregate effort. Our anonymity proposition ---
that $\Delta$ depends on the entrant count alone and is invariant to the
resource distribution --- is a neutrality statement of the same shape, and it is
precisely what makes P9 tractable. Whether the two are formally related, or
merely rhyme, is not established here and should not be asserted without proof.

\section{Hopkins and Kornienko: the standard result, read in full}
\label{sec:hk}

\citet{HopkinsKornienko2010} is the paper Referee~2 invoked when predicting that
the standard result --- inequality reduces status expenditure --- would survive
in this setting. It has now been read in full \textbf{[F]}, and the reading
changes the positioning in four specific ways.

\paragraph{Primitives.} A continuum of contestants with endowment $z$ drawn from
$G$, private information. Each divides $z$ between visible performance $x$ and
private consumption $y = z - x$. A continuum of indivisible rewards with values
distributed $H$ is assigned assortatively by rank in performance, so that in a
separating equilibrium $G(z) = F(x) = H(s) = r$. Utility is $U(x, y, s)$ with
monotonicity, complementarity and own-concavity conditions. The unique symmetric
separating equilibrium solves a differential equation with the non-competitive
optimum as boundary condition, and equilibrium performance strictly exceeds the
cooperative level everywhere above rank zero, so the equilibrium is Pareto
dominated from the contestants' point of view. Existence and uniqueness are
obtained by fitting the model into \citet{Mailath1987}.

\paragraph{Point 1: there is no extensive margin.} Every contestant competes.
The decision is how much performance to buy, never whether to compete at all.
No participation constraint, no entry cost, no fixed fee. The binary entry
decision that generates $k^*$ in P5 has no counterpart in the paper. This is the
cleanest available statement of what the rebuilt model adds and it is now
verified rather than inferred.

\paragraph{Point 2: the contest is perfectly discriminating.} The authors note
explicitly that much of the surrounding contest literature assumes a partly
stochastic award mechanism, and describe their own model by contrast as a
perfectly discriminating rank-order tournament. The rebuilt model is on the
stochastic side: scores are $\mu r_i + (1-\mu) s_i e_i$ and the induced CDFs
$F$ and $G$ carry all the content. The two models sit in different branches, and
P-MU --- where the sign of $\partial \Delta(0) / \partial Q$ turns on a
hump-shaped weighting --- is an object that cannot arise without noise.

\paragraph{Point 3: the wealth channel is a budget, not a fee.} Endowment enters
as the resource divided between performance and consumption. It sets how much
performance a contestant can afford at the margin. It does not act through the
utility cost of a fixed payment, which is the entire content of
$\kappa(w) = u(w) - u(w-c)$ in the rebuild. These are different mechanisms and
should not be described as the same one.

\paragraph{Point 4 --- the most useful finding: the right stochastic order is
already named.} The comparative statics are conducted in the \emph{dispersive
order} and the \emph{star order}, not in mean-preserving spreads. Their
Definition~1, taken from \citet{ShakedShanthikumar2007}, is that $F$ is less
dispersed than $G$, written $F \leq_d G$, whenever $G^{-1}(r) - F^{-1}(r)$ is
weakly increasing on $(0,1)$; equivalently, at a fixed rank the more dispersed
distribution is the less dense one. They also record that there is no clear
relation between the dispersive order and second-order stochastic dominance,
because SOSD compares distributions on both mean and dispersion while the
dispersive order concerns dispersion alone.

This bears on open item~2 of \texttt{PROOFS.tex}, which observes that
Theorem~\textsf{P9gen} requires $T$ to be a pivot-spread and that a general
Rothschild--Stiglitz spread need not be monotone, leaving the general case open.
The dispersive order is exactly the tool that literature uses for this
difficulty, and the two conditions are close. Writing $T = \Lambda_\beta^{-1}
\circ \Lambda_\alpha$, the dispersive order says $T(w) - w$ is increasing; if in
addition $T(w) - w$ changes sign, it does so once, at a point $x_0$, giving
$T(w) \geq w$ above $x_0$ and $T(w) \leq w$ below it --- which is the
pivot-spread hypothesis. So the pivot-spread class is not an ad hoc restriction
invented to make the proof work; it is, up to the crossing condition, the
standard dispersion order of this literature. That is a considerably stronger
position than the current text claims for itself, and it should be checked
formally and then stated.

\paragraph{Direction of their result, stated carefully.} Under rank indexing,
their Proposition~4 assumes the minimum endowment is higher ex post, endowments
are less dispersed ex post, and the maximum endowment is lower ex post. The
conclusion is that performance is \emph{higher} ex post on $[0, \hat r]$, where
$\hat r$ is the single crossing point of the two endowment functions, and that
utility rises at the bottom but falls for the middle and top. So less dispersion
raises competition at the bottom and middle. Read in the other direction, more
dispersion of endowments lowers performance there --- which is the ``standard
result'' Referee~2 had in mind, and it does hold.

Two cautions follow for P9. First, their conclusion is stated on an interval
below a crossing point, not globally, so it is not a uniform-sign aggregate
claim and should not be described as one. Second, their model already contains a
crossing point at which effects change sign, though it is a crossing in the
\emph{endowment functions} with a corresponding sign change in \emph{utility},
not a sign change in the direction of the expenditure response. P9's pivot rule
is a different object: a sign reversal in $\mathrm{d}k^*/\mathrm{d}(\text{spread})$
located at the pivot $x_0$ relative to the marginal entrant's wealth
$w_{(k^*)}$. The claim in \texttt{PROOFS.tex} \S11 that the result
``reconciles'' the standard result with the motivating pattern is defensible
only if stated at that level of precision. As drafted it is too loose, because
it invites the reading that the standard result is uniform in sign, which the
source does not say.

\paragraph{Genealogy of the dispersive order, corrected.} \citet{HopkinsKornienko2004}
and \citet{HopkinsKornienko2009} have now been read in full \textbf{[F]}, and
together they change what was said about the dispersive order in an earlier
draft of this section. The claim there was that the order traces to the general
stochastic-orders textbook of \citet{ShakedShanthikumar2007} and merely happens
to fit this problem. That understates it. The dispersive order, applied to
exactly this class of problem, is Hopkins and Kornienko's own methodological
innovation, developed across these two papers before the 2010 paper uses it.

\citet{HopkinsKornienko2004} is the earlier, income-indexing treatment: agents
are indexed by income $z$ itself, and equilibrium strategies and utilities are
written as functions of $z$. Comparative statics there use first- and
second-order stochastic dominance, together with two likelihood-ratio
refinements --- the monotone likelihood ratio (MLR) order and, for spreads, the
unimodal likelihood ratio (ULR) order. Under the ULR order the paper's
Proposition~4 partitions the population into three zones by two threshold
points $\hat z_-, \hat z_+$ --- poor, middle, and rich --- and signs the effect of
a spread separately in each zone, with the middle zone (where the density ratio
is highest) unambiguously seeing higher conspicuous consumption. This is a
\emph{two}-crossing-point result, more general than a single pivot.

\citet{HopkinsKornienko2009} is explicit, in its own introduction, about why a
second methodology was needed: the income-indexing approach of the 2004 paper
cannot compare distributions with different supports, so it cannot handle a
policy that raises the floor of the income distribution while taxing the top.
The paper switches to indexing agents by \emph{rank} $r = G(z)$ instead, and
introduces the dispersive order for this purpose, citing \citet{Shaked1982}
directly (full citation added to \texttt{refs.bib} below) --- rank-indexing
plus the dispersive order is presented there as a new technique built
specifically to sign comparative statics of a competitive equilibrium under a
change in the underlying distribution, which is precisely the class of problem
P9-gen belongs to.

Structurally, \citet{HopkinsKornienko2009} Propositions~4 and~5 are close
cousins of P9-gen. Proposition~4 gives a sufficient condition for a full Pareto
improvement: the poorest agent no worse off, together with the distribution
becoming \emph{more} dispersed in their sense, is enough for utility to rise
almost everywhere. Proposition~5 is the mirror image: the poorest richer, the
richest poorer, and the distribution \emph{less} dispersed, then utility falls
for the middle and upper ranks --- this is the source proposition the 2010
paper's Proposition~4 (discussed above in \S\ref{sec:hk}) generalises. In both
cases the sign of the welfare effect is pinned by a single crossing point
$\hat r$ of the two distribution functions, in exactly the way P9-gen pins the
sign of the entry-count effect by the pivot $x_0$. The two constructions are not
the same theorem --- theirs crosses in the distribution functions and signs a
welfare comparison, P9-gen crosses in $T(w) - w$ and signs an entry-count
comparison --- but they are the same technique applied to structurally
analogous problems, and P9-gen should cite this lineage rather than only the
general textbook.

\paragraph{A concrete open question this raises.} \citet{HopkinsKornienko2004}'s
ULR result permits \emph{two} crossing points and three sign regions, not one.
P9-gen's pivot-spread hypothesis is a single-crossing object. Whether a
two-pivot generalisation of P9-gen exists --- for spreads that compress agents
in the middle of the wealth distribution while still expanding the tails, which
a single pivot cannot represent --- is an open question raised directly by this
reading and is not addressed anywhere in the current bundle.

\paragraph{Consequence for the income/wealth question.} Both papers work
entirely in \emph{income} space: the primitive is $z$, described throughout as
income, divided in a single period between a positional good and a
non-positional good, with no saving and no mention of wealth as a stock. This
is independent confirmation, from the second major reference paper in this
literature (after Schroyen--Treich), of the point made when the primitive in
\texttt{PROOFS.tex} was discussed: a one-period resource divided between an
entry cost and residual consumption is standard practice as ``income'' in this
literature, not ``wealth,'' and the field's own vocabulary supports the
relabelling already adopted.

\paragraph{No entry margin, confirmed again.} As with Hopkins--Kornienko
(2010) and Schroyen--Treich, there is no participation decision in either
paper: $x \in [0, z/p]$ is a continuous choice on a convex feasible set. The
extensive-margin gap the entry\_contest bundle fills is confirmed a third time.

\citet{Hopkins2012} and \citet{Hopkins2008} are adjacent, not yet read in full.
\textbf{[C]}

\section{Endogenous entry into contests, read in full}
\label{sec:entry}

Four papers in this strand have now been read in full \textbf{[F]}:
\citet{LevinSmith1994} and \citet{MorganOrzenSefton2012} in their published
versions, \citet{FuJiaoLu2015} in its published version, and \citet{FuLu2010}
as MPRA Paper 945 (November 2006), the working version of the \emph{Economic
Inquiry} article. Together they settle the entry-structure question more
sharply than the abstracts allowed.

\paragraph{The identity problem, stated by the literature itself.} All four
papers assume identical potential entrants. Every one of them then confronts the
same difficulty: with identical agents, the pure-strategy count is determined
but the identity of who enters is not. \citet{LevinSmith1994}, p.~586 (main
text; footnote~2 is attached to the preceding sentence and concerns
Engelbrecht-Wiggans), say of the prior literature's deterministic asymmetric
equilibria that the process by which potential bidders divide into entrants and
non-entrants is not explained. \citet{MorganOrzenSefton2012}, p.~441, say in their own model that
because all players are identical the identity of those opting in is not
uniquely determined. The literature has three devices for this. Levin and Smith
restore symmetry by mixed entry, so the number of entrants becomes stochastic
(binomial). \citet{FuLu2010} (footnote~6) and \citet{MorganOrzenSefton2012}
(Proposition~1) use sequential entry with observability, so arrival order breaks
the tie; Morgan et al.\ add private random delays to make the order well
defined. \citet{FuJiaoLu2015} take the mixed route and, because entrants do not
observe the number of rivals, need Dasgupta--Maskin's theorem for
two-dimensional discontinuous games to establish existence (their Theorem~1).

P5 is the fourth device, and it differs from the other three in kind: the
resource ordering $\kappa(w_{(1)}) < \kappa(w_{(2)}) < \cdots$ is a property of
the \emph{agents}, not of a timing protocol or a randomisation, so the
pure-strategy count $k^*$ is deterministic under \emph{simultaneous} entry with
no mixing and no arrival order --- and, by the count-invariance proposition in
\texttt{PROOFS.tex}, the same in every pure-strategy equilibrium.

The claim should stop there, and an earlier version of this survey did not.
Heterogeneous costs do not by themselves pin \emph{which} agents enter: with
$\Delta=(12,10,1)$ and $\kappa=(2,5,8)$ the sets $\{1,2\}$, $\{1,3\}$ and
$\{2,3\}$ are all equilibria, one of them with the richest challenger absent.
What is delivered is the count, plus existence of the assortative equilibrium,
plus its uniqueness under a stated extra condition; absent that condition,
assortativity is a selection. That is still a real advance on mixing, arrival
order and random delays, each of which leaves even the assortative outcome
undefined --- but it is a smaller claim than ``identity is a property of the
agents'', and the two footnotes cited above are the places to hang the smaller
one.

\paragraph{Levin--Smith.} $N$ identical risk-neutral potential bidders, fixed
entry cost $c$ paid ex ante before learning one's value (contrast Samuelson's
interim costs), $n$ revealed before bidding. Their $n^*$ is defined as the
largest integer at which the marginal entrant's expected gain is non-negative
--- the same object as $k^*$ for identical agents --- and their symmetric
equilibrium is in mixed strategies with $n \sim \mathrm{Binomial}(N, q^*)$.
Two results bear on the open welfare item (R2 Comment~3). In common-value
auctions free entry is \emph{excessive} by a business-stealing argument
(Proposition~3, citing \citet{MankiwWhinston1986}), and the seller taxes entry;
in independent-private-value auctions free entry is \emph{optimal} because the
marginal entrant's private gain equals the social gain, via $V_n - W_n = n(V_n
- V_{n-1})$ (Proposition~6, equation~18).

Which branch applies is settled by a \emph{criterion}, not by resemblance:
Proposition~6 holds exactly when equation~(18) holds, and (18) requires the
marginal entrant's private gain to equal the increment to total value. With a
prize fixed at $V$ we have $V_n \equiv V$, so $V_n - V_{n-1} = 0$: the social
gain from any entrant beyond the first is \emph{exactly zero} while the social
cost is $c>0$. The paper states this directly (p.~596): ``In CV auctions,
social gains are zero (and therefore smaller than social costs) for all
$n \geq 2$.'' The entry\_contest model is therefore in the common-value branch
\emph{by construction} --- $V$ exogenous and fixed \emph{is} $V_n \equiv V$ ---
and the benchmark is Proposition~3, with the expected verdict of excessive
entry. Only a prize that grew with participation could reinstate
Proposition~6.

Their Proposition~9 adds that welfare falls monotonically as the pool $N$ grows
beyond $n^*$, through the coordination cost of mixed entry; Proposition~8 is
the supporting statement that the probability of \emph{no} entry rises with $N$
under an optimal mechanism. That is a different mechanism from P7's cap, which
is a count bound with no coordination content, but both say that a larger pool
does not help beyond a point.

Their equation~(9), $(1-q)^{N-1} V = c$, requires care. It is \emph{not} an
entry-equilibrium condition: it is the first-order condition of the welfare
problem, obtained by differentiating their equation~(8),
$S = [1-(1-q)^N]V - qNc$, and setting $\partial S/\partial q = 0$. It
characterises the socially optimal entry probability, which the seller's
optimal entry fee $e^*$ is chosen to induce; the paper's own proof of
Proposition~8 writes it as $(1-q^s_N)^{N-1} = c/V$, with $s$ for social. The
bidders' equilibrium condition is their equation~(6), $B_i(q^*,\Omega)=0$.
Equation~(9) is nonetheless algebraically identical to \citet{FuJiaoLu2015}'s
Definition~1, which reaches the same equation from the opposite direction --- as
a \emph{lower bound} on equilibrium entry, derived from the win-only-if-alone
deviation payoff. The coincidence is worth stating, but the two must not be
identified: in common-value auctions the free-entry equilibrium lies
\emph{above} $q^s$ (Proposition~3), and reading (9) as an equilibrium condition
erases the very wedge that Proposition~3 is about.

\paragraph{Fu--Jiao--Lu, and the accounting behind P7.} $M \geq 2$ identical
risk-neutral potential bidders, common fixed entry cost, generalised nested
lottery contest, prize purse capped at $V$ and allowed to be contingent on the
realised number of entrants. Their equation~(2) is a rent-accounting
inequality: expected prize mass paid out is at least expected entry costs plus
expected bidding costs, so
\[
  [1-(1-q)^M]\,V \;\geq\; Mq\,(\Delta + \mathbb{E}[x^\alpha]),
\]
which implies that the expected number of entrants $Mq$ cannot exceed
$V/\Delta$, since the bracket is at most $1$ and $\mathbb{E}[x^\alpha] \geq 0$.
Cite \emph{Definition~2 alongside equation~(2)} for this: Fu et al.\ never
display the count bound, and what they derive from (2) is a chain about
\emph{bids} --- (3), then (4) by convexity, (5), the bid bound (6), and
Theorem~2's single-peakedness. The count bound is instead the content of their
Definition~2, $\bar q = \arg\max\{q \in [0,1] : (1-(1-q)^M)V \geq Mq\Delta\}$,
every feasible point of which satisfies $Mq\Delta \leq V$. A referee sent to
equation~(2) alone will find only the bid chain and conclude the attribution is
loose.
That is the same accounting logic as P7. P7's content is therefore not a new
inequality but the deterministic, heterogeneous, utility-unit version of a
known one: it bounds a pure-strategy count rather than an expected count, it
uses $\kappa$ in utility units rather than a nominal cost, and it holds under
fixed rules with no designer. The paper should say exactly that. What Fu et
al.\ build on top --- the uniform upper bound on expected overall \emph{bids}
(Theorem~2), its single-peakedness in $q$, its implementation by Tullock
contests with entry fees (Theorem~7), and the optimal shortlisting rule
(Lemma~6, Theorem~9, with $\bar M = \min\{N : V/N < \alpha\Delta/(\alpha-1)\}$)
--- are design results with an effort-maximising organiser and have no
counterpart in the entry\_contest model, which has no designer and a single
fixed prize. Their Theorem~11, that the optimal entry probability, the bid
bound and the shortlist size all fall in the entry cost, is the design-side
cousin of P8.

\paragraph{Fu--Lu.} $M \geq 3$ identical, sequential entry with full
observation of current participants, organiser chooses prize and a per-entrant
fee or subsidy. Their equation~(7), total effort equals budget minus $N$ times
the entry cost, is again the P7 accounting identity, and is what drives their
Theorem~1 that the optimal contest attracts exactly two entrants. Their
conclusion lists different types of contestants as future research. So the
entire design literature on endogenous entry --- Fu--Lu 2010, Fu--Jiao--Lu 2015
--- is homogeneous-agent by construction.

\paragraph{Morgan--Orzen--Sefton, and what the experiment did not vary.} The
theory section gives $n^* = \lfloor \sqrt{P/F} \rfloor$ for identical agents
with equal endowment $w$, outside option $F$ and prize $P$. The experiment holds
$w = 100$ for every subject by design. Predicted entry was~2 (small prize)
and~4 (large prize); observed was 2.5 and 3.7 in the second half (rounds
26--50; 2.7 and 3.6 in rounds 1--25). The \emph{signs differ}: the small prize
shows \emph{excess} entry (2.5 against 2, significant in both halves,
$p = 0.004$ and $0.012$) while the large prize shows a \emph{shortfall} (3.7
against 4, differing from Nash only at the 6\% and 12.5\% levels). Both
deviations run against what within-contest behaviour would predict --- with
negative returns to two-person competition one expects \emph{fewer} than two,
and with under-investment one expects \emph{more} than four --- and the paper
says of each case ``we observe precisely the opposite''. Citing this as
``observed entry exceeds prediction'' would be wrong for half the experiment.
Earnings were \emph{not}
equalised across the contest and the outside option in either treatment, and a
self-selection-by-risk-attitude explanation is rejected by comparing
three-player contests across treatments (small-prize triples over-invest and
earn less than the outside option, large-prize triples under-invest and earn
more; the earnings difference is significant at 1\%). The same controlled
comparison also defeats a ``contest premium'' explanation, so two behavioural
accounts are ruled out, not one. Two things follow. First, the experimental
literature on contest entry has not tested resource sorting, because endowments
were held equal; a Morgan--Orzen--Sefton design with heterogeneous endowments
is the natural test of P8--P9. Second, their concluding observation that the
marginal entrant's power is blunted because inframarginal rent-seeking moves
the marginal calculus has no analogue in the binary model, where inframarginal
behaviour is fixed at ``invest''; that is why the cutoff is clean here and
noisy there.

\paragraph{Moreno--Wooders, read in full, and the relabelling objection answered.}
\citet{MorenoWooders2011} has now been read in its published form \textbf{[F]}.
$N$ risk-neutral buyers in a standard independent-private-value auction; each
buyer's entry cost $Z_i$ is an independent draw from $H$ on $[\underline c,
\bar c]$ and is \emph{privately} observed before entry; the value is learned
after entry; the seller may set a screening value $v$ and an admission fee
$\phi$. Their Proposition~2 gives a unique symmetric equilibrium in threshold
form: enter iff $z < t^*(v,\phi)$, where $t^*$ solves $U(v, H(t)) = t + \phi$
and is decreasing in both instruments. The number of bidders is
$\mathrm{Binomial}(N, H(t^*))$. Proposition~3: $v = \phi = 0$ maximises social
surplus, because the marginal type's private gain equals its social
contribution. Proposition~4: inframarginal entrants keep information rents, so
the seller no longer captures the whole surplus. Propositions~5--8: the
revenue-maximising screening value is positive but below the fixed-$n$ level;
with an admission fee available, screen by entry cost rather than by value; an
entry cap plus fee does better still. Proposition~9: heterogeneity is
irrelevant asymptotically --- only the lower bound $\underline c$ matters.
Section~6: revenue may rise \emph{or} fall in $N$, unlike Levin--Smith. There
are no comparative statics in the distribution $H$ anywhere in the paper.

The relabelling objection --- is $\kappa \circ \Lambda^{-1}$ substantively
different from a primitive cost distribution? --- can now be answered on four
points, and the first exposes a gap in \texttt{PROOFS.tex}.

\emph{Information.} Moreno--Wooders' costs are private, so the equilibrium is
a symmetric Bayesian threshold and only the \emph{distribution} of the entrant
count is pinned. \texttt{PROOFS.tex} writes the entry condition as
$\kappa(w_{(j)}) \leq \Delta(j-1)$ in terms of the order statistics of the
wealth profile, which presupposes that the profile is common knowledge; that is
what makes $k^*$ deterministic and pins the identities of the entrants. This
assumption is nowhere stated in \texttt{PROOFS.tex}. It must be, because it is
the first structural difference from Moreno--Wooders and the reason P5 delivers
an exact count where they deliver a binomial.

\emph{Primitive.} Their $H$ is a primitive and they never vary it. Ours is
induced from a resource distribution through a known concave map, so a
mean-preserving spread of $\Lambda$ is a specific, non-arbitrary deformation
of the induced cost distribution. P9 has no counterpart in Moreno--Wooders
because there is no distributional comparative static there to have a
counterpart to.

\emph{Designer.} Theirs is mechanism design with three instruments; ours has
none.

\emph{Objective.} Their Proposition~3 is the independent-private-value logic
where the marginal entrant's private gain equals the social gain. For a
fixed-prize winner-take-all contest the common-value logic of
\citet{LevinSmith1994}'s Proposition~3 applies instead, and free entry is
excessive. Their Proposition~4 carries over directly: entrants with
$w > \bar w$ keep surplus.

So the threshold-in-cost-space is theirs; the count-pinning under complete
information, the induced-cost structure, and the distributional comparative
static are not. The objection is answered --- but only once the information
assumption is written down.

\paragraph{Still at abstract level.} \citet{KaplanSela2010} give a rationale
for entry fees in an all-pay auction with privately known entry costs.
\textbf{[C]} \citet{Ginzburg2021} has a continuum of contestants and a
profit-maximising entry price. \textbf{[A]} \citet{GrossmannDietl2015} has
high- and low-talent contestants choosing between entry and heterogeneous
outside options. \textbf{[A]} \citet{Samuelson1985} is the interim-entry-cost
variant (cost incurred after learning one's value) that both Levin--Smith and
Moreno--Wooders position against. \textbf{[C]}

\section{Budget and liquidity constraints}

\citet{CheGale1997} has now been read in full \textbf{[F]}, in its published
form. The setting is $N \geq 2$ risk-neutral bidders with common value $v$,
publicly known wealths $w_1 > \cdots > w_N$, and a hard cap: a bidder cannot
spend more than his wealth (no borrowing in \S3; credit at rate $r$ in \S4).
Utility is linear; wealth enters only as the cap. There is no entry cost and no
participation decision.

Their Proposition~1 (all-pay auction) gives expected total expenditure
$\min\{v, w_2\}$ with only the two wealthiest active --- the rich preempt the
poor. Their Lemmas~1--3 (lottery) give a unique equilibrium with a critical
bidder $i^*$: those with wealth above a common bid $b^*$ all bid $b^*$, those
below bid their wealth, and everyone participates. Proposition~3: the lottery
dissipates more than the all-pay auction iff $v$ exceeds a threshold above
$w_2$. Adding low-wealth entrants has no effect in the all-pay auction and
always raises lottery expenditure.

This reading confirms first-hand what \S\ref{sec:st} previously took from
Schroyen--Treich's description: a budget constraint is a kink in an otherwise
linear payoff at $w_i$, and nothing else. The entry\_contest model has no kink
and no cap; $\kappa(w)$ is smooth and the sorting comes from curvature
everywhere. The critical-bidder structure in their lottery is the
hard-constraint cousin of the $\kappa$ threshold, but it operates on the
intensive margin (bid size) while ours operates on the extensive margin (enter
or not), and in their lottery participation is never a choice. Their all-pay
cap $\min\{v, w_2\}$ is wealth-driven and depends on the second-highest
endowment; P7's cap $V/\kappa$ is prize-driven and independent of the
distribution. Their observation that adding low-wealth agents leaves all-pay
expenditure unchanged rhymes with anonymity plus P7 --- agents who would not
enter do not move $k^*$ --- and is worth a sentence.

\citet{GrossmannDietl2012} consider two bidders with asymmetric valuations under
imperfect capital markets and find aggregate investment lower when at least one
is constrained. \textbf{[A]} \citet{GrossmannHottiger2020} and
\citet{CheGale1998} continue the line. \textbf{[C]}

\section{Risk aversion and concavity in contests, read in full}

\citet{CornesHartley2012} has been read in full as Manchester Economics
Discussion Paper EDP-0806 (May 2008), the working version of the
\emph{Economic Theory} article. \textbf{[F, WP version.]} The setting is $n
\geq 2$ contestants with logistic CSF and concave production, an indivisible
common rent $R$, continuous expenditure, concave utility $u_i$, complete
information; no entry cost.

The multiplicity example is now concrete. Their Example~1 has ten identical
contestants with $u(c) = c - 0.45c^2$ --- increasing absolute risk aversion ---
and possesses both the symmetric equilibrium with each spending $0.0563$ and an
asymmetric one in which three contestants spend $0.184$ and seven spend zero;
the active three earn more than in the symmetric equilibrium. Uniqueness
(Theorem~4.2) needs \emph{regularity}, and their sufficient condition
(Lemma~4.2) is a curvature bound $2u'(R-x) \geq u'(-x)$ on $(0,R)$; it holds
for CRRA when initial wealth $I_i \geq R(1 - 2^{-1/\gamma})^{-1}$
(Corollary~4.1), for small rents (Corollary~4.2), for prudent contestants with
absolute risk aversion at most $1/(2R)$ (Corollary~4.3), and for CARA. Whether
non-increasing absolute risk aversion suffices is left open.

The contrast with P5 is now exact. Their asymmetric equilibrium is a
coordination among \emph{identical} agents on who is active, sustainable
because the continuous aggregative best response admits several fixed points
under IARA. P5 has agents ordered by $\kappa(w)$, so who is active is pinned
by the order and there is nothing to coordinate on; P5 needs only that
$\kappa$ is decreasing, with no curvature bound of the Lemma~4.2 type. That is
the reason the binary-entry structure escapes their multiplicity, and it should
be stated in one sentence with their Example~1 as the foil.

Two further results are directly relevant. Their dropout value $\bar Y_i =
f_i'(0)[u_i(R) - u_i(0)]/u_i'(0)$ (Lemma~3.2) makes a contestant inactive
whenever aggregate lobbying exceeds it --- a curvature threshold on the
\emph{aggregate}, where $\kappa(w)$ is a curvature threshold on \emph{own}
resources. And Proposition~6.2 is a selection result: in large asymmetric
contests only the types with the highest dropout value stay active, so
competition drives the more risk-averse out and the limiting dissipation ratio
$H[u] = [u(R) - u(0)]/(Ru'(0))$ is that of the least risk-averse type. That is
selection by risk attitude across heterogeneous $u_i$ with common wealth; ours
is selection by resource level with common $u$. The two are complementary and
should be cited together. Their CRRA specification $u(c) = (I_i + c)^{1-\gamma}
/(1-\gamma)$ has initial wealth as an additive shift in a static game, the same
one-period object as in Schroyen--Treich.

\citet{Treich2010} has been read as LERNA working paper 09-013 of 17 February
2009 \textbf{[F, WP version]}. $n \geq 2$ identical agents with initial wealth
$w$, concave $u$, rent $b$, payoff $p_i u(w + b - x_i) + (1-p_i) u(w - x_i)$,
general CSF. Proposition~2: prudence ($u''' \geq 0$) is sufficient for risk
aversion to lower effort relative to risk neutrality, for any $n$ and any
regular CSF; Corollary~1: for $n = 2$ it is necessary and sufficient, and
quadratic $u$ gives no effect at all. Proposition~1: DARA plus a small rent
gives a unique symmetric equilibrium. The paper's closing sentence lists entry
as future work. Two things follow. The risk-aversion strand needs $u'''$ to
sign anything, whereas P8 and P9-gen need only $u''$; that is a concrete
statement of what ``no higher-order restrictions'' buys. And wealth here is
again a one-period additive shift. \citet{SkaperdasGan1995} is the antecedent;
\citet{CornesHartley2005} is the general-technology companion;
\citet{CornesHartley2003} establishes regularity under CARA. \textbf{[C]}

\section{Competitor numbers and the shape of noise, read in full}

\citet{RyvkinDrugov2020} has been read in its published \emph{Theoretical
Economics} form \textbf{[F]}, and \citet{DrugovRyvkin2020} as the working
paper of 13 March 2019 \textbf{[F, WP version]}.

The setting is symmetric risk-neutral players, output $y_i = e_i + X_i$ with
i.i.d.\ additive noise, a winner-take-all prize normalised to one, strictly
convex cost, and a number of players $K$ that may be deterministic or
stochastic; heterogeneity is assumed away explicitly (their footnote~3). There
is no entry margin. The two objects that drive everything are
\[
  b_k = \mathbb{E}\big[f(X_{(k-1:k-1)})\big], \qquad
  E^*_k = k\,c'^{-1}(b_k) = \mathbb{E}\big[h(X_{(k-1:k)})\big]
  \text{ under quadratic cost},
\]
the density of noise at the largest of the other $k-1$ draws, and the hazard
rate at the second-largest of $k$. Proposition~1 and Corollary~1: a unimodal
noise density is inherited by individual effort as a function of $k$, by a
Karlin variation-diminishing argument on $F^{k-1}(1-F)$. Corollary~4: monotone
density gives monotone effort, constant iff uniform. Proposition~2: symmetric
unimodal density gives $e^*_2 = e^*_3$ and decreasing thereafter.
Proposition~3: IFR noise with cost more convex than quadratic gives aggregate
effort increasing in $K$; DFR noise with cost less convex than quadratic gives
aggregate effort decreasing. Proposition~5: for a mean-preserving increase in
uncertainty about $K$, log-concave noise lowers aggregate effort and log-convex
raises it. Proposition~7 (Appendix~A.1) gives existence and uniqueness of the
symmetric pure-strategy equilibrium under a lower bound on $c''$ and bounds on
$f$ and $f'$.

The relationship to P-MU can now be stated at the level of objects rather than
by analogy. P-MU's $\Delta(0)$ is the win-probability gain to the first
investor, a functional of the induced score distributions; its dependence on $Q$
runs through the same order-statistic weighting that gives $b_k$ its dependence
on $k$. Proposition~1's inheritance of unimodality is the structural template
for P-MU's hump-shaped weighting, and the sentence on their p.~1589 --- that no
prediction about the effect of competitor numbers survives for every noise
distribution --- is exactly the frame in which R2's $Q$-conjecture should be
refuted. The next step is mechanical: write $\Delta(0)$ in the form
$\int \phi \, dF_{(\cdot)}$ and check whether their log-supermodularity
condition on $-\tilde G_\theta$ is what P-MU's weighting needs. Their footnote~29
records that \citet{DrugovRyvkin2020a} shows the dispersive order is necessary
and sufficient for ranking equilibrium effort across noise distributions --- the
dispersive order again, now in the contest branch as well as the status branch.

The 2019 working paper is the prize-allocation companion. With prizes summing
to one and monotone in rank, the optimum is a two-prize schedule with the top
$r^*$ ranks paid equally (Proposition~1), where $r^*$ maximises $B_{r,n}/r =
\tfrac{1}{n}\mathbb{E}[h(Z_{(n-r:n)})]$; IFR gives winner-take-all, DFR gives
reward-all-but-last (Proposition~2); the convex transform order gives monotone
comparative statics in skewness (Proposition~4); log-concave noise makes
effort Schur-convex in the prize vector (Proposition~7). Proposition~8 applies
this to status tournaments in the sense of \citet{MoldovanuSelaShi2007}: with
IFR noise the top status category is a singleton, with DFR it contains $n-1$.
This is not P-MU territory --- the entry\_contest model has one prize --- but it
is the status-tournament link back to the abandoned paper's motivation, and a
referee from that literature will expect it cited. \citet{LimMatros2009} is the
binomial-$K$ special case. \textbf{[C]}

\section{Assignment models and the pivot rule, read in full}
\label{sec:assignment}

Two papers previously listed as ``assignment-model comparators'' turn out to
carry the closest published precedent for P9-gen's pivot rule.
\citet{CostrellLoury2004} has been read as the draft of 11 December 2003
\textbf{[F, draft version]} and \citet{Suen2007} in its published form
\textbf{[F]}.

\paragraph{Costrell--Loury and the single-pivot precedent.} In their two-job
model (\S2), jobs are filled in proportions $\theta$ and $1-\theta$, and the
entire wage schedule is determined by one sufficient statistic: $\hat\mu =
F^{-1}(\theta)$, the ability of the \emph{marginal} worker between jobs. Any
shift that raises $\hat\mu$ raises low-ability wages and lowers high-ability
wages. Under a mean-preserving spread on fixed support, they state that the
effect depends on whether the marginal worker at quantile $\theta$ lies in the
upper or lower tail: if $\theta$ is high the spread raises $\hat\mu$ and
narrows the wage distribution, if $\theta$ is low it lowers $\hat\mu$ and
widens it. That is a single-pivot sign rule located at the marginal agent's
quantile relative to the spread's crossing point --- the same logic as P9-gen,
where the sign of $\mathrm{d}k^*$ flips according to whether $w_{(k^*)}$ lies
above or below $x_0$. Their object is a wage schedule through $\hat\mu$; ours
is an entry count through $w_{(k^*)}$. P9-gen should be presented as the
entry-count analogue of the Costrell--Loury two-job comparative static. Their
continuum result (Proposition~6) then shows what governs the direction when
there is no single marginal agent: with $G$ riskier than $F$ on $[0,1]$,
concave $\beta$ lowers $w(0)$ and widens the span, convex $\beta$ lowers $w(1)$
and narrows it. The curvature of the payoff map decides. For us that map is
$\kappa$.

Three further results are tools. Their Lemma~1 (in the proof of
Proposition~5) shows that $G$ riskier than $F$ on fixed support implies the
quantile functions reverse the order, $F^{-1}$ riskier than $G^{-1}$; their
Proposition~5 signs an aggregate (output) under a general spread with multiple
crossings by integrating $\Gamma(\rho) = \int_\rho^1 [G^{-1} - F^{-1}]$ against
$\mathrm{d}\beta$; and their Theorem~1 shows that a Blackwell-more-informative
signal induces a second-order-riskier distribution of posterior means, so any
mean-preserving-spread result has an information reading. Their Proposition~10
(Cobb--Douglas crowding) removes the curvature dependence altogether: a spread
always narrows the span. The draft is not the published article; the
published version should be checked for these proposition numbers.

\paragraph{Suen and the general-spread technique.} One-to-one transferable-
utility matching with $Q = \theta(x)\phi(y)$, positive assortative matching, and
matching function $\mu = G^{-1} \circ F$, which they note can itself be read as
a distribution function. Proposition~2: if $G_1$ is more dispersed than $G_0$
(same mean, second-order), then total earnings on the other side fall, under
concavity of $\theta, \phi$ and log-concavity of $1-F$. The proof handles a
\emph{general} Rothschild--Stiglitz spread with an arbitrary finite number of
crossings $k_0 < k_1 < \cdots < k_{2n}$ of the quantile functions: the
quantile difference alternates in sign, each partial integral is signed by
second-order dominance, concavity controls the recombination, and the result
is integrated by parts against $\rho = (1-F)/f$. Proposition~3 and its
corollaries give conditions under which every agent loses. Suen's $\mu$ is
cross-side (the other market's distribution moves) whereas our $T =
\Lambda_\beta^{-1} \circ \Lambda_\alpha$ is same-side (before and after), but
the map class is identical.

\paragraph{Consequence for open item~2.} Two published sources sign an
aggregate under a general mean-preserving spread, by the same first move but
under \emph{different} hypotheses --- and they are not independent: Suen's
proof imports the quantile-reversal step from Costrell--Loury by citation
(``see, for example, [2]''), and his footnote~1 engages their rank-rescaling
device directly. One lineage, not two convergent confirmations.

Both begin by decomposing the quantile difference at the crossing points.
They then diverge:

\begin{itemize}
\item \citet{CostrellLoury2004}'s Proposition~5 integrates
  $\Gamma(\rho) = \int_\rho^1 [G^{-1} - F^{-1}]$ against $\mathrm{d}\beta$ and
  needs only that $\beta$ is \textbf{non-decreasing}. It requires no concavity
  and no hazard-rate condition; the strings ``log-concav'' and ``hazard'' do
  not occur anywhere in that paper. Concavity of $\beta$ belongs to their
  Proposition~6, a different result about wage \emph{inequality}.
\item \citet{Suen2007}'s Proposition~2 integrates by parts against
  $\rho = (1-F)/f$ and needs concavity of $\theta,\phi$ \emph{and}
  log-concavity of $1-F$ --- the latter being exactly equivalent to
  $\rho' \leq 0$, since $\mathrm{d}^2 \log(1-F)/\mathrm{d}x^2 = \rho'/\rho^2$.
\end{itemize}

The weaker hypothesis is Costrell--Loury's, so the \textbf{monotone-weight
route should be tried first}, not log-concavity.

Three obstacles stand between this and a result, none of them previously
recorded here.

\begin{enumerate}
\item \textbf{Aggregate versus count.} Both sources sign a smooth integral
  ($Q$, $W$). $k^*$ is an integer count defined by a threshold crossing.
  Neither source signs a count, and this is the substantive gap.
\item \textbf{The rescaling trap.} The obvious route writes $k^*$ as a
  functional of $\kappa \circ \Lambda^{-1}$ --- a rank-space rescaling, which
  is precisely what Suen's footnote~1 declines to do: ``unless the density
  function $g$ is increasing, concavity of $\phi$ does not imply concavity of
  $\hat\phi$''. Curvature does not survive the change of variable. Since the
  dispersive-order machinery of \S\ref{sec:hk} that motivates
  $T = \Lambda_\beta^{-1}\circ\Lambda_\alpha$ \emph{is} rank-space machinery,
  this sits on the main road, not a side path. Either establish curvature in
  the rescaled variable directly, or accept an increasing-density restriction
  on $\Lambda$ that the model does not currently make.
\item \textbf{Cross-side versus same-side.} Suen compares two populations
  across a market; we compare one population before and after a spread. The
  map class is identical; the object compared is not.
\end{enumerate}

This remains a concrete, checkable lead. It has not been tried and should not
be claimed until it has.

\section{Positional competition and education}

\citet{ChungLee2017} is the nearest thing to the abandoned paper's framing:
inequalities in wages, productivity and education affordability determining
educational effort in competition for premier jobs, in a Tullock framework with
endogenous returns, finding that competition for better-paying jobs can produce
socially inefficient excessive effort and that a widening wage gap intensifies
the competition. \textbf{[A]} Two players, continuous effort.

\section{Foundations for a rank-allocated prize, read in full}
\label{sec:foundations}

Three papers previously cited secondhand have now been read in full, and they
do a specific job: they justify the premise that \texttt{PROOFS.tex} currently
asserts, that a fixed prize $V$ is allocated by rank rather than by price.

\paragraph{Cole--Mailath--Postlewaite 1992.} Published form \textbf{[F]}.
Multigenerational; a continuum of men with capital and women with a
non-traded endowment; joint consumption; CRRA. Status is defined as a ranking
device that determines how one fares in the non-market sector. Two social
norms: wealth-is-status and aristocratic (inherited). Proposition~1: a
wealth-is-status equilibrium exists (Mas-Colell's continuum-game theorem).
Property~1: in equilibrium no family's relative rank ever changes, even though
everyone competes. Property~2: everyone weakly over-saves relative to the
no-matching benchmark, except the bottom agent, who is undistorted.
Proposition~2: an aristocratic equilibrium exists if the capital distribution
is sufficiently spread out at the lower tail. Their footnote~20 says a market
in which relative position determines the house one buys \emph{cannot}
generate the multiplicity they are after --- the prizes must be non-market.
Their \S IV.A two-period example states that the agent in the economy with the
more compact income distribution tends to have the higher savings rate, and
\S V.D that once the distribution is dispersed enough the incentive disappears.
That is the 1992 origin of the ``standard result'' R2 invoked, on the intensive
margin.

\paragraph{Cole--Mailath--Postlewaite 1995.} CARESS Working Paper 95-14,
1 August 1995, marked forthcoming in the Minneapolis Fed \emph{Quarterly
Review} \textbf{[F, WP version]}. The working-paper title has ``Concern''
singular; Hopkins--Kornienko's 2010 reference list has ``Concerns'' plural at
\emph{QR} 19(3), 12--21. The published title has not been confirmed and the
bib entry carries the working-paper form. The concern for relative standing is
\emph{instrumental}: rank matters only because a non-market decision depends
on it. \S 2 is a complete-information effort model with the matching function
$m(y)$ equal to the distribution function of output, so $m'$ is large when the
distribution is tight and competition is intense (\S 2.1, \S 2.3). \S 3 is the
incomplete-information version: wealth is unobservable, so the rich signal by
destroying wealth, and the fraction destroyed rises with rank. Footnote~8:
it is competition from \emph{below} that distorts effort --- truncate the
population from the top and nothing changes; truncate from the bottom and the
new lowest agent becomes undistorted. \S 4: when prizes are allocated by rank
rather than sold, the welfare theorems do not apply.

\paragraph{What these two give the entry\_contest model.} Finding~5 asks why
threshold entry into a rank-allocated prize is not already covered. The
answer to the ``rank-allocated prize'' half is that Cole--Mailath--Postlewaite
is where that premise is justified and where its consequences (no welfare
theorems, non-market allocation, instrumental status) are worked out.
\texttt{PROOFS.tex} should cite them at the point where $V$ is introduced and
say that $V$ is a prize in their sense --- not sold, allocated by rank. Their
``competition from below'' is the same structural fact as the anonymity
proposition and the ordering by $\kappa$: only agents who could outrank you
matter. Their Property~1 is consistent with the model here --- entry does not
change the wealth rank; it is the score, not the rank, that the prize follows.
And their multiplicity is \emph{across} norms, not within one, so it is not in
tension with P5.

\paragraph{Becker--Murphy--Werning.} \citet{BeckerMurphyWerning2005}, draft revised December 2003 \textbf{[F,
draft version]}. Utility $u(c,s)$ with the crucial assumption $u_{cs} > 0$:
status raises the marginal utility of consumption. A fixed supply of status
positions, sold in a hedonic market (or, equivalently by their Proposition~1,
assigned by rank in consumption of a fixed-supply, intrinsically worthless
social good). Stage one: fair lotteries over income. Proposition~2: if the
constant-marginal-utility allocation is a mean-preserving spread of the initial
income distribution, it is the competitive lottery equilibrium --- a minimum
equilibrium inequality that any more compact initial distribution is spread
out to reach. Proposition~3: with a status market this coincides with the
utilitarian planner's allocation, so the planner \emph{increases} inequality.
\S 7: when status attaches to income \emph{rank} rather than being purchased,
rank lotteries impose real externalities (one's win lowers others' rank) and
the equilibrium no longer coincides with the planner's. This is not a contest
paper; the mechanism is orthogonal to $\kappa$, the mean-preserving spread is
an \emph{output} of equilibrium rather than an input, and $u_{cs} > 0$ is the
structural opposite of our separable $V$. It matters for one thing only: the
welfare section (R2 Comment~3). Their \S 7 states in print that rank-allocated
status carries real externalities that purchased status does not, which is the
zero-sum, business-stealing frame the welfare benchmark needs, and their
Proposition~3 is a warning that a planner may \emph{want} inequality.

\citet{ChungLee2017}, \citet{ColeMailathPostlewaite1998},
\citet{Postlewaite1998}, \citet{Frank1985}, \citet{FernandezGali1999} and
\citet{BilanciniBoncinelli2012} remain at abstract or citing-paper level.

\section{Refuted conjectures in context}
\label{sec:refuted}

Referee~2 conjectured that for $\mu > 0$, higher $Q$ shrinks the pie among
investors so fewer invest. \texttt{PROOFS.tex} \S R2 refutes this: the sign
depends on $\mu$, and the crossover $\mu^*$ is distribution-dependent. Placed
against \citet{RyvkinDrugov2020}, this refutation is the expected outcome rather
than a surprise, and framing it that way is both more honest and more
persuasive than presenting it as an isolated counterexample.

\section{Background and toolkit}

\paragraph{Lazear--Rosen, read in full, and a precedent that was hiding in
plain sight.} \citet{LazearRosen1981}, published form \textbf{[F]}. The
risk-neutral two-player tournament: $C'(\mu) = (W_1 - W_2) g(0)$, efficient,
with the prizes in equation~(10) equal to expected product plus or minus
$V/2g(0)$, which they read as an \emph{entrance fee or bond}: each player posts
a bond and takes a fair winner-take-all gamble over the pool. Footnote~2: a
pure-strategy equilibrium exists only when the noise variance is large enough
--- the ancestor of Ryvkin--Drugov's Proposition~7. \S III adds risk aversion,
and the subsection titled ``Income Distributions'' contains the result that
matters here. With a common utility function and different \emph{endowed}
incomes $y_0$, decreasing absolute risk aversion makes the rich prefer the
tournament and the poor prefer the piece rate, so agents \emph{self-select into
the contest by wealth}, and the resulting income distribution is positively
skewed. That is the first appearance, in 1981, of ``who enters a contest is
determined by wealth''. The channel is risk attitude --- DARA, so the wealthy
find the gamble cheaper --- not the utility cost of a fixed fee, and it needs
$u'''$ where $\kappa$ needs only $u''$. But it is a precedent for the
mechanism and it should be cited as one, with the channel distinction stated.
\S IV (heterogeneous ability, asymmetric information) gives adverse selection
into the top league and a competitive handicap $h^* = \Delta\mu/2$; the
handicap algebra is the tool for R2's open question about an equilibrium in
which the incumbent abstains and a challenger invests.

\paragraph{Moldovanu--Sela and Hoppe--Moldovanu--Sela, read in full.}
\citet{MoldovanuSela2001} as the draft of 5 June 2000 \textbf{[F, draft
version]}: $k \geq p$ contestants with \emph{private ability} entering a
separable cost $c_i \gamma(x)$, all-pay, perfectly discriminating; a single
first prize is optimal under linear or concave cost (Propositions~3.2, 4.2)
and two prizes can be optimal under convex cost (Proposition~4.3). Their \S 5
entry fee produces a threshold in \emph{ability} space --- a third threshold
variant, alongside Moreno--Wooders' cost space and our resource space. Their
footnote~12 warns that independence of types is problematic with endogenous
entry, which applies here too: the entrant pool is the top $k^*$ by $w$, not an
i.i.d.\ sample. \citet{HoppeMoldovanuSela2009}, published form \textbf{[F]}:
finite two-sided matching on costly signals with private attributes, built on
\citet{BarlowProschan1966}'s normalised-spacings results. Proposition~4: entry
on the long side raises own-side signalling for every $F$, and moves the
other side's signalling up or down according to whether $F$ is DFR or IFR;
Example~1 shows welfare falling in $n$ up to $n = 8$ in a small market and
rising thereafter. Proposition~5: the relatively \emph{more homogeneous} side
signals more, because it perceives a more dispersed prize structure. That is
the perfectly-discriminating, private-ability counterpart of the
competitor-number comparative static: with no noise, entry always raises
own-side effort; P-MU says that with noise the sign at the margin is
distribution-dependent even for symmetric agents. Their IFR/DFR spacings
machinery and Ryvkin--Drugov's hazard-rate machinery are the same
Barlow--Proschan lineage.

\paragraph{Surveys and standard references.} \citet{Konrad2009} and
\citet{Corchon2007} are the two surveys a referee will expect to see
acknowledged; the 2010 Hopkins--Kornienko text cites Konrad for the point that
greater heterogeneity among competitors and a smaller spread of prizes both
reduce equilibrium effort. \textbf{[F, as characterised there.]}
\citet{MoldovanuSela2006} is the contest-architecture companion and
\citet{SchotterWeigelt1992} the experimental treatment of asymmetric
tournaments. \textbf{[C]}

Two references carry technical weight rather than positioning weight.
\citet{ShakedShanthikumar2007} is the source for the dispersive and star orders
and is the right citation for the pivot-spread class discussed in \S3.
\citet{RothschildStiglitz1970} is the definition open item~2 of
\texttt{PROOFS.tex} appeals to when it observes that a general mean-preserving
spread admits only a mean-preserving coupling $w' = w + \varepsilon$ with
$\mathbb{E}[\varepsilon \mid w] = 0$, which need not be monotone. Citing both,
and stating that the pivot-spread hypothesis is the dispersive-order
restriction rather than an arbitrary one, converts that open item from an
admitted weakness into a deliberate and named choice of stochastic order.

\section{Verdict on novelty}

The mechanism alone does not carry a paper, and after the full read of
\citet{SchroyenTreich2016} it carries even less than the previous draft of this
document supposed. ``Resource-sorted entry via concavity'' is their privilege
contest plus a \citet{MorenoWooders2011} threshold on the entry structure, and
the resource-label distinction does not separate the two (\S\ref{sec:st}).
Nor is a mean-preserving spread in the resource distribution new in itself:
their Theorem~3 already signs one inside a privilege contest.

Two further sharpenings from the full reads. \citet{LazearRosen1981}'s
``Income Distributions'' subsection (within \S III) already has agents
self-selecting into a tournament by endowed wealth, through DARA; the mechanism
has a 1981 precedent and the contribution sentence must acknowledge it. Two
qualifications, both verified, bound how much it concedes.

First, \emph{their claim is an example and they say so}. The subsection opens
``While it is not possible to make a general argument based on an example,
table~1 suggests\ldots'', and rests on one utility ($U = ay^a$), quadratic
costs, normal errors and two endowment levels. The honest acknowledgement is
that Lazear and Rosen \emph{suggest} the sorting on an example under DARA,
while the present model \emph{derives} it under concavity alone,
distribution-free.

Second, the channel difference ($u''$ here, $u'''$ there) is substantive rather
than verbal, and separates decisively. Our $\kappa(w) = u(w)-u(w-c)$ is
strictly decreasing whenever $u'' < 0$, with no condition on $u'''$; DARA needs
$A'(w) < 0$, whose numerator is $(u'')^2 - u'''u'$, so it \emph{requires}
$u''' > 0$. Quadratic utility witnesses the gap: it is concave, has $u''' = 0$,
and is \emph{increasing} absolute risk aversion --- our sorting operates on it,
theirs cannot.

A third difference, which the channel formulation understates, is the
\emph{object of choice}: theirs is a choice between two compensation schemes
that both pay, so the sorting is tolerance for the binomial spread of a
lottery; ours is whether to pay a fixed fee at all, so the sorting is
affordability in utility units. That is different economics, not merely a
different derivative, and it is the more robust half of the answer.

And \citet{MorenoWooders2011}
(\S\ref{sec:entry}) shows that what separates P5 from a private-cost threshold
is the complete-information assumption on the wealth profile --- which
\texttt{PROOFS.tex} uses but does not state.

What is not in any source located here is the combination of three results that
require the $Q$-player binary-entry setting jointly:

\begin{enumerate}
\item \textbf{P7}, the uniform cap $k^* \leq V/\kappa$ holding for every $Q$
  with no limit taken. The accounting logic behind it --- entrants times entry
  cost cannot exceed prize mass --- is already in \citet{FuJiaoLu2015},
  equation~(2) \emph{together with their Definition~2} (they never display the
  count bound itself; see \S\ref{sec:entry}), and \citet{FuLu2010},
  equation~(7), for identical agents with mixed or sequential entry. Fu--Lu's
  is the sharper case, (7) being an equality rather than an inequality --- but
  their $N$ is the count in an \emph{optimally designed} contest, so ``under
  fixed rules with no designer'' is doing real work. P7's content is the
  deterministic heterogeneous version in utility units, bounding a
  pure-strategy count rather than an expected one, under fixed rules with no
  designer. It should be presented as that, not as a new inequality.
\item \textbf{P9 and P9-gen}, the pivot rule locating the sign reversal of the
  inequality effect on aggregate entry at $w_{(k^*)}$ versus $x_0$. The novelty
  is not that a mean-preserving spread is signed --- \citet{SchroyenTreich2016}
  do that --- but that it is signed \emph{globally and without higher-order
  restrictions on $u$}, where the intensive-margin result is local and needs
  the CSF decisiveness parameter and the third derivative. The literature
  supplies the stochastic order (dispersive) but not the entry-count
  application. The closest structural precedent for the pivot rule itself is
  \citet{CostrellLoury2004}'s two-job model (\S\ref{sec:assignment}), where the
  sign of a spread's effect flips according to whether the marginal worker's
  quantile lies above or below the crossing; P9-gen is its entry-count
  analogue and should cite it as such.
\item \textbf{P-MU}, the exact condition for the sign of the competitor-number
  effect in a binary-entry contest, as the discrete analogue of the
  \citet{RyvkinDrugov2020} \emph{density} result
  $b_k = \mathbb{E}[f(X_{(k-1:k-1)})]$ --- \emph{not} their hazard-rate result,
  which governs \emph{aggregate} effort as
  $\mathbb{E}[h(X_{(k-1:k)})]$ and uses a different order statistic (the two
  densities stand in ratio $k(1-F)$). Since $\Delta(0)$ is an individual gain,
  $b_k$ is the right comparator. Note that \citet{SchroyenTreich2016} list an
  arbitrary number of players as an extension they did not undertake, so the
  $Q$-dimension is unoccupied on both sides.

  The comparison has now been computed rather than conjectured, and it is
  exact: P-MU's weight is $W = G^{Q-1}(1-G)F$, whose hump-shaped factor
  $G^{Q-1}(1-G)$ peaks at $G = (Q-1)/Q$, while Ryvkin and Drugov's
  log-supermodular kernel $-H_\theta = F^{k-1}(1-F)$ peaks at $F = (k-1)/k$.
  Under $G \leftrightarrow F$, $Q \leftrightarrow k$ these are the \emph{same
  function}; our extra factor $F$ is the incumbent's CDF and carries no $Q$.
  So the hump-shape is not an artefact of our construction but the standard
  object of this literature. Log-supermodularity transfers (the cross-partial
  of $\log[z^{n-1}(1-z)]$ is $1/z > 0$), but the second hypothesis does not:
  Karlin's argument needs the integrand in the role of $u'$ to be single
  crossing $+-$, whereas ours is $F\cdot(f-g) = -F\varphi'$, which crosses
  $-+$. \emph{No unimodality claim for $\Delta(0,Q)$ in $Q$ is available};
  the plausible target of a reworked argument is an interior \emph{minimum},
  which would sit consistently with R1--R2, and it has not been derived.
\end{enumerate}

The contribution should be stated as those three theorems, not as the mechanism.
That is a narrower claim than the rebuild handover assumed, and it survives.

\section{What still needs doing}

Read in full as of this version: Hopkins--Kornienko 2004, 2009, 2010;
Schroyen--Treich (WP); Ryvkin--Drugov (TE, published); Drugov--Ryvkin (JET,
WP); Fu--Jiao--Lu (published); Fu--Lu (WP); Cornes--Hartley (WP); Treich
(WP); Levin--Smith (published); Morgan--Orzen--Sefton (published);
Moreno--Wooders (published); Che--Gale 1997 (published); Suen (published);
Costrell--Loury (draft); Cole--Mailath--Postlewaite 1992 (published) and
1995 (WP); Becker--Murphy--Werning (draft); Moldovanu--Sela 2001 (draft);
Hoppe--Moldovanu--Sela (published); Lazear--Rosen (published). Every
substantive claim in this document about what a paper does now rests on a full
read, except where marked [A] or [C]. Where a working paper or draft was read,
the published version has not been line-checked.

Outstanding, in descending order of consequence for the current claims:

\begin{enumerate}
\item State the complete-information assumption on the wealth profile in
  \texttt{PROOFS.tex} (\S\ref{sec:entry}). Without it P5's deterministic
  $k^*$ is not distinguished from Moreno--Wooders' binomial count.
\item Add \citet{LazearRosen1981} to the contribution paragraph as the
  precedent for wealth-based self-selection, with the $u''$ versus $u'''$
  distinction.
\item Cite \citet{ColeMailathPostlewaite1992,ColeMailathPostlewaite1995} at
  the point where $V$ is introduced (\S\ref{sec:foundations}).
\item Relate P-MU's weighting $W$ to \citet{RyvkinDrugov2020}'s $b_k$ and
  their log-supermodularity condition --- a computation.
\item Try the Suen / Costrell--Loury general-spread technique on P9-gen
  (\S\ref{sec:assignment}) --- a proof attempt.
\item Published-version checks: Schroyen--Treich (\emph{GEB}, Theorem~3 and
  the payoff specifications); Fu--Lu (\emph{EI}, the title differs between
  MPRA and Fu--Jiao--Lu's citation); Cole--Mailath--Postlewaite 1995
  (\emph{QR}, ``Concern'' or ``Concerns''); Costrell--Loury,
  Becker--Murphy--Werning, Moldovanu--Sela, Cornes--Hartley, Treich,
  Drugov--Ryvkin (proposition numbering).
\item \citet{GrossmannDietl2015} and \citet{ChungLee2017}, for the
  entry-condition comparison and the welfare template respectively --- still
  at abstract level.
\end{enumerate}

\bibliographystyle{plainnat}
\bibliography{refs}

\end{document}
